../cictikz/src/cictikz/data/tex/cictikz_preamble.tex

% Every transition a two-bit binary code can make: four states, and all
% twelve ordered pairs between them. The point of the slide is the count —
% twelve ways to get it wrong — so every arrow is drawn, including both
% directions of each pair.

\begin{tikzpicture}[thick]

  ../cictikz/src/cictikz/data/tex/dacsm_lib.tex

  \dacState{s00}{(0,3)}{00}
  \dacState{s01}{(3.6,3)}{01}
  \dacState{s11}{(0,0)}{11}
  \dacState{s10}{(3.6,0)}{10}

  % The two directions of a pair are drawn as separate arrows, offset either
  % side of the line between the states, as the original does.
  \dacEdge{s00}{s01}   \dacEdge{s01}{s00}
  \dacEdge{s11}{s10}   \dacEdge{s10}{s11}
  \dacEdge{s00}{s11}   \dacEdge{s11}{s00}
  \dacEdge{s01}{s10}   \dacEdge{s10}{s01}
  \dacEdge{s00}{s10}   \dacEdge{s10}{s00}
  \dacEdge{s01}{s11}   \dacEdge{s11}{s01}

\end{tikzpicture}
\end{document}
